
\documentclass{article}
\usepackage{colortbl}
\usepackage{makecell}
\usepackage{multirow}
\usepackage{supertabular}

\begin{document}

\newcounter{utterance}

\centering \large Interaction Transcript for game `codenames', experiment `abstract', episode 0 with reflexion{-}Qwen3.5{-}27B.
\vspace{24pt}

{ \footnotesize  \setcounter{utterance}{1}
\setlength{\tabcolsep}{0pt}
\begin{supertabular}{c@{$\;$}|p{.15\linewidth}@{}p{.15\linewidth}p{.15\linewidth}p{.15\linewidth}p{.15\linewidth}p{.15\linewidth}}
    \# & $\;$A & \multicolumn{4}{c}{Game Master} & $\;\:$B\\
    \hline

    \theutterance \stepcounter{utterance}  
    & & & \multicolumn{2}{p{0.3\linewidth}}{
        \cellcolor[rgb]{0.9,0.9,0.9}{
            \makecell[{{p{\linewidth}}}]{
                \texttt{\tiny{[GM$|$GM]}}
                {'hidden': {'team': ['adventurousness', 'interconnection', 'oppressiveness', 'anticommunism'], 'innocent': ['fantasy', 'endearment', 'ethic', 'solemnity', 'adversity', 'manner', 'loathsomeness'], 'opponent': [], 'assassin': []}, 'revealed': {'team': {'team': [], 'innocent': [], 'opponent': [], 'assassin': []}, 'opponent': {'team': [], 'innocent': [], 'opponent': [], 'assassin': []}}}
            }
        }
    }
    & & \\ \\

    \theutterance \stepcounter{utterance}  
    & & \multicolumn{4}{p{0.6\linewidth}}{
        \cellcolor[rgb]{0.9,0.9,0.9}{
            \makecell[{{p{\linewidth}}}]{
                \texttt{\tiny{[P1$\langle$GM]}}
                \texttt{You have attempted this same task before and did not succeed. These are the plans you wrote after those attempts.} \\
\texttt{Trial \#1: New plan: My previous strategy failed because I prioritized covering many words with broad clues like "state" and "belief," which were too ambiguous and led my partner to guess an opponent's word ("unification"). Next time, I will prioritize safety over speed by selecting clues that are highly specific to a small subset of my words (1{-}2 words) to minimize the risk of triggering opponent or distractor words. I will avoid abstract concepts that could easily overlap with the opponent's list (e.g., "unification" vs. "interconnection") and instead use concrete, unique associations. For example, instead of "belief" for "affirmation" and "anticommunism," I might use "yes" for "affirmation" alone, or "cold" for "anticommunism" alone, ensuring the partner has a clear, narrow path to the correct answer without guessing related but incorrect words.} \\
\\ 
\texttt{Trial \#2: New plan: My previous strategy failed because I prioritized safety by giving clues for only one word at a time, which was too slow and allowed the opponent to catch up or the game to end before I cleared my list. Specifically, I missed the opportunity to group "adventurousness" and "interconnection" (or "oppressiveness") under a broader but safe clue like "journey" or "system," and I failed to group "anticommunism" with "oppressiveness" under a clue like "regime" or "control." Next time, I will aim for a balance by grouping 2{-}3 words per clue using specific semantic links that do not overlap with the opponent's remaining words (e.g., "grandiosity," "viciousness"). I will avoid single{-}word clues unless the remaining words are completely unrelated, and I will explicitly check if a broader clue (like "travel" for adventurousness/interconnection) risks triggering an opponent word before committing.} \\
\\ 
\texttt{Trial \#3: New plan: My previous attempt failed because I selected a clue ("system") that was too broad and abstract, likely causing my partner to guess a distractor word like "ethic" or "manner" instead of the specific targets. I also failed to utilize "adventurousness" and "anticommunism," leaving half my list unaddressed. Next time, I will prioritize high{-}specificity clues that create a unique semantic bridge to multiple words while strictly avoiding overlap with the distractor list. Specifically, I will use "risk" to link "adventurousness" and "anticommunism" (as a political risk), and "control" to link "oppressiveness" and "interconnection" (as a controlling network), ensuring the clues are concrete enough to prevent guessing "fantasy" or "solemnity." I will also ensure I attempt to cover all four words in the first turn if a safe, unifying clue exists, rather than splitting them into disjointed, low{-}yield groups.} \\
\\ 
\\ 
\texttt{Provide one single{-}word clue that relates to some of your team's words from the list below. You can choose to how many and to which words your clue relates to.} \\
\texttt{The goal is to have your partner guess all of your team's words based on the clues you give them, before the other team has uncovered all of theirs. If your partner guesses an 'opponent' or a 'distractor' word, their guessing turn ends and the opposing team gets to make a turn. If your partner guesses a 'kill' word, the game ends and your team loses immediately.} \\
\texttt{The more words your clue relates to, the quicker you can win the game, but the harder the associations may be to guess for your partner, so choose your clue word wisely. The clue word has to be semantically related to the target words, it cannot be one of the words in the lists or contain parts of them.} \\
\\ 
\texttt{Always give your single{-}word clue and your comma{-}separated list of related target words in the following format and make your answers as short as possible, never include any other text than is required in this form:} \\
\\ 
\texttt{CLUE: <WORD>} \\
\texttt{TARGETS: <TARGETS>} \\
\\ 
\texttt{Your team words are: adventurousness, interconnection, oppressiveness, anticommunism.} \\
\texttt{Your opponent's team words are: .} \\
\texttt{Distractor words are: fantasy, endearment, ethic, solemnity, adversity, manner, loathsomeness.} \\
\texttt{The kill words are: .} \\
            }
        }
    }
    & \\ \\

    \theutterance \stepcounter{utterance}  
    & \multicolumn{4}{p{0.6\linewidth}}{
        \cellcolor[rgb]{0.9,0.9,0.9}{
            \makecell[{{p{\linewidth}}}]{
                \texttt{\tiny{[P1$\rangle$GM]}}
                \texttt{CLUE: control} \\
\texttt{TARGETS: oppressiveness, interconnection} \\
            }
        }
    }
    & & \\ \\

    \theutterance \stepcounter{utterance}  
    & & & \multicolumn{2}{p{0.3\linewidth}}{
        \cellcolor[rgb]{0.9,0.9,0.9}{
            \makecell[{{p{\linewidth}}}]{
                \texttt{\tiny{[GM$|$GM]}}
                \texttt{control} \\
            }
        }
    }
    & & \\ \\

    \theutterance \stepcounter{utterance}  
    & & & \multicolumn{2}{p{0.3\linewidth}}{
        \cellcolor[rgb]{0.9,0.9,0.9}{
            \makecell[{{p{\linewidth}}}]{
                \texttt{\tiny{[GM$|$GM]}}
                ['oppressiveness', 'interconnection']
            }
        }
    }
    & & \\ \\

    \theutterance \stepcounter{utterance}  
    & & & \multicolumn{2}{p{0.3\linewidth}}{
        \cellcolor[rgb]{0.9,0.9,0.9}{
            \makecell[{{p{\linewidth}}}]{
                \texttt{\tiny{[GM$|$GM]}}
                {'word': 'oppressiveness', 'assignment': 'team'}
            }
        }
    }
    & & \\ \\

    \theutterance \stepcounter{utterance}  
    & & & \multicolumn{2}{p{0.3\linewidth}}{
        \cellcolor[rgb]{0.9,0.9,0.9}{
            \makecell[{{p{\linewidth}}}]{
                \texttt{\tiny{[GM$|$GM]}}
                {'word': 'interconnection', 'assignment': 'team'}
            }
        }
    }
    & & \\ \\

    \theutterance \stepcounter{utterance}  
    & & & \multicolumn{2}{p{0.3\linewidth}}{
        \cellcolor[rgb]{0.9,0.9,0.9}{
            \makecell[{{p{\linewidth}}}]{
                \texttt{\tiny{[GM$|$GM]}}
                \texttt{opponent has won} \\
            }
        }
    }
    & & \\ \\

    \theutterance \stepcounter{utterance}  
    & \multicolumn{4}{p{0.6\linewidth}}{
        \cellcolor[rgb]{0.95,0.95,0.95}{
            \makecell[{{p{\linewidth}}}]{
                \texttt{\tiny{[P2$\rightarrow$P1]}}
                \texttt{Post{-}episode Reflexion (not sent to GM):} \\
\texttt{New plan: My previous attempt failed because I selected "control" as a clue, which was too broad and likely triggered the distractor "manner" or "ethic," causing a loss. Additionally, I only targeted two words ("oppressiveness", "interconnection"), leaving "adventurousness" and "anticommunism" unaddressed, which is inefficient. Next time, I will prioritize a clue that specifically links "adventurousness" and "anticommunism" (e.g., "risk" or "danger") while ensuring it does not semantically overlap with the distractors like "adversity" or "fantasy." I will also aim to group all four words under a single, highly specific abstract concept if possible, or split them into two very tight pairs with distinct, non{-}overlapping clues to avoid triggering distractors. I will explicitly verify that my chosen clue does not share a strong semantic link with any word in the distractor list before submitting.} \\
            }
        }
    }
    & & \\ \\

\end{supertabular}
}

\end{document}
